% ID: SEL-D0030-OB-E03
% ESTADO: cotejada. Segunda pasada visual completa y borrador compilado revisados.
% Fuente: Selectividad/matematicas_aplicadas/titular_junio/Examen-Matemáticas Aplicadas a las Ciencias Sociales.pdf
% SHA-256: 5bd7838829b9f653746319907f86642017d12e66ed097db0126fb544a2d0f632
% Página: 2; opción B; ejercicio original 3.

Un campus universitario dispone de 3000 plazas numeradas de aparcamiento para vehículos, distribuidas en tres zonas A, B y C. La zona A está constituida por las plazas del 1 al 1500, estando 1350 de ellas protegidas del sol. La zona B la conforman las plazas numeradas desde 1501 a 2500, estando el 80\% protegidas del sol. La zona C contiene las plazas numeradas desde 2501 hasta 3000, estando solamente 250 protegidas del sol. Aleatoriamente se elige una de las plazas de aparcamiento del campus.
\begin{enumerate}[label=\alph*)]
\item \textbf{(0.75 puntos)} \textquestiondown{}Cuál es la probabilidad de que esté en la zona A o en la B?
\item \textbf{(0.75 puntos)} \textquestiondown{}Cuál es la probabilidad de que no esté protegida del sol?
\item \textbf{(1 punto)} Si se ha elegido una plaza protegida del sol, \textquestiondown{}cuál es la probabilidad de que esté ubicada en la zona B?
\end{enumerate}
